\section*{Chapter \ref{ch:iv:fixed-income-and-markets} --- Fixed Income and Markets}

\begin{solution}{iq:iv:fixed-income-and-markets:1}
$-7 \times 0.25\% = -1.75\%$, slightly less in magnitude once convexity is counted.
\iqlookfor{duration times the move, with the sign.}
\end{solution}

\begin{solution}{iq:iv:fixed-income-and-markets:2}
$10\,000\,000 \times 0.98 \times 6.5 \times 0.0001 = 6\,370$ per basis point.
\iqlookfor{DV01 from price and modified duration, with the price applied.}
\end{solution}

\begin{solution}{iq:iv:fixed-income-and-markets:3}
$1.04^2/1.03 - 1 \approx 5.01\%$: a rising curve implies forwards above the spot rates.
\iqlookfor{the no-arbitrage relation between zero and forward rates.}
\end{solution}

\begin{solution}{iq:iv:fixed-income-and-markets:4}
Driver, rates move by maturity and the curve, cross-asset links with reasons, what surprised you, what you watch next, all in
two minutes with rounded numbers (\cref{met:iv:fixed-income-and-markets:recap}). The interviewer does not want a list of closing
levels, a move with no cause, or confident causes for moves that had none.
\iqlookfor{structure and causality over recall.}
\end{solution}

\begin{solution}{iq:iv:fixed-income-and-markets:5}
Match DV01s: $10 \times 850/460 \approx 18.5$ million of the second bond. The hedge removes parallel-move risk at the two
maturities' average; curve risk (the two points moving differently), basis risk between the two issues, and convexity remain.
\iqlookfor{DV01-neutral sizing and the residual risks.}
\end{solution}

\begin{solution}{iq:iv:fixed-income-and-markets:6}
$D_{\mathrm{mod}} \approx 8.11$ and $\mathcal C \approx 80.8$. Duration alone: $-8.11\%$. With convexity: $-8.11\% + 0.40\% =
-7.71\%$. Exact: $-7.72\%$ (the bond priced at 5\%). Convexity recovers almost all of the error of the linear estimate.
\iqlookfor{the second-order expansion and its accuracy.}
\end{solution}

\begin{solution}{iq:iv:fixed-income-and-markets:7}
The five-year yield is 3.69\%. Carry: the coupon of 3.69 minus 3\% repo on a price of 100, about 0.69\%. Roll-down: repriced at
the four-year yield of 3.59\%, the bond gains about 0.33\%. Total about 1.02\% if the curve does not move; with a duration near
3.7 at the horizon, a rise of about 28 basis points in the four-year yield would wipe it out.
\iqlookfor{separating carry from roll-down and computing the break-even move.}
\end{solution}

\begin{solution}{iq:iv:fixed-income-and-markets:8}
$(4\% - 4.5\%) \times 100\,000\,000/12 \approx -41\,667$: negative carry. The trade needs the bond's price to rise by more than
that (about 4 cents per 100 of price over the month), from falling yields or a richening of the bond against the curve.
\iqlookfor{the sign of carry and what the position is betting on.}
\end{solution}

\begin{solution}{iq:iv:fixed-income-and-markets:9}
$1.10 \times 1.04/1.02 \approx 1.1216$. A market forward of 1.1180 is below parity: buying euros forward is cheap relative to
the money-market route, which is the cross-currency basis (dollars are expensive to obtain through FX swaps), not a free lunch
for most participants because of balance-sheet costs.
\iqlookfor{covered interest parity and the basis as the explanation of deviations.}
\end{solution}

\begin{solution}{iq:iv:fixed-income-and-markets:10}
A swap needs no funding of the notional and little balance sheet, whereas holding a government bond ties up balance sheet and
repo capacity; when dealers' balance sheets are costly or supply of long bonds is heavy, bonds cheapen relative to swaps and the
spread turns negative. Earning it means buying the bond and paying fixed on the swap, financing the bond in repo: the trade
earns the spread only if repo stays cheap and it can be held to maturity, and it is exposed to mark-to-market losses if the
spread widens further (One Quant Book~2, chapter~9).
\iqlookfor{balance-sheet cost as the explanation and the funding risk of the trade.}
\end{solution}

\begin{solution}{iq:iv:fixed-income-and-markets:11}
``A payroll surprise pushed front-end rate expectations up: two-year yields $+12$, ten-year $+5$, a 7 basis point flattening.
The dollar rallied 0.8\% on the wider rate differential. Equities fell on the rate shock and recovered half by the close as
investors weighed stronger growth. Credit spreads did not move, which says the market read it as a growth story, not a stress
one.'' A surprise would have been a long end rising more than the front end, which would suggest term premium or inflation
concerns rather than the policy path.
\iqlookfor{a causal recap that uses the numbers and flags what would not fit.}
\end{solution}

\begin{solution}{iq:iv:fixed-income-and-markets:12}
The duration-matched barbell holds about 59.6\% in the two-year and 40.4\% in the thirty-year; its convexity is about 173
against about 81 for the ten-year bullet. On a parallel move of 100 basis points either way the barbell does better, which the
chapter's code checks. The catch: the extra convexity is priced, since the barbell usually yields less, so it loses if yields do
not move, and the curve can move non-parallel (a steepening hurts the barbell's long leg).
\iqlookfor{duration matching, the convexity comparison and the cost of convexity.}
\end{solution}

\begin{solution}{iq:iv:fixed-income-and-markets:13}
Front end up, since the path of policy rates is repriced higher. Long end up less, possibly down if the market reads the hike
as lowering future inflation (a flattening, even an inversion). Currency up on the rate differential, unless the hike is read as
a sign of an inflation problem or as damaging growth. Equities down on higher discount rates, though financials may gain and
the market may rally if uncertainty is removed. Each ``other way'' is a statement about what the move says for growth, inflation
and credibility, and a good answer names which reading it is assuming.
\iqlookfor{causal reasoning across assets with explicit assumptions.}
\end{solution}
